% animations.tex — SwiftUI + UIKit animations: implicit vs explicit, Animatable,
% transitions, matchedGeometryEffect, springs, Phase/KeyframeAnimator,
% Transaction; UIView.animate, UIViewPropertyAnimator, Auto Layout, model vs
% presentation.
% Sources (checked 2026-10-06):
%   developer.apple.com/documentation/swiftui/animation/default (a spring since iOS 17;
%     easeInOut before) · /swiftui/view/animation(_:value:) (iOS 13) · the unconstrained
%     /swiftui/view/animation(_:) (iOS 13, deprecated 15: "Use withAnimation or
%     animation(_:value:)") · /swiftui/view/transaction(_:) · /transaction/disablesanimations ·
%     /swiftui/animatablepair · /swiftui/view/matchedgeometryeffect (iOS 14) ·
%     /swiftui/phaseanimator · /keyframeanimator · /keyframetrack (iOS 17)
%   developer.apple.com/documentation/uikit/uiview/animate(springDuration:bounce:…) (iOS 17) ·
%     /uikit/uiviewpropertyanimator (iOS 10) · /continueanimation(withtimingparameters:…) ·
%     /uiviewcontrollercontexttransitioning/completetransition(_:) ·
%     /quartzcore/calayer/presentation()
% @labels: area=swiftui kind=api platform=ios topic=ui discussion=172
% @tags: withanimation, implicit-animation, transaction, animatabledata, transitions, matchedgeometryeffect, spring-animation, phaseanimator, keyframeanimator, uiview-animate, uiviewpropertyanimator, presentation-layer
\documentclass[8pt]{extarticle}
\usepackage{printup-sheet}

\lstdefinelanguage{SwiftSheet}{
  morekeywords={struct,class,final,func,var,let,some,if,else,return,in,self,
    true,false,nil,get,set,private,static,View,Shape,CGFloat,Bool},
  sensitive=true, morecomment=[l]{//}, morestring=[b]"}

\tikzset{
  lbl/.style={font=\scriptsize, text=black!75, inner sep=1pt},
  st/.style={box, minimum height=6mm, font=\scriptsize, inner sep=2pt},
}

\begin{document}

\sheettitle{Animations — SwiftUI \& UIKit}{swiftui · folio}

\oneliner{You never animate frame by frame: you \textbf{change state} (SwiftUI)
or set the \textbf{final value} (UIKit) \emph{inside an animation context}, and
the framework interpolates every \textbf{animatable} property old $\to$ new.
The questions are always: \emph{which} change is in the context, and \emph{can}
it be interpolated?}

\begin{multicols}{2}
\raggedright

\section{SwiftUI}
\begin{itemize}
  \item \textbf{Explicit} \texttt{withAnimation(.spring) \{ open.toggle() \}}:
        animates \emph{every} view depending on state changed in the block
        (mutation site decides). \textbf{Implicit}
        \texttt{.animation(.easeIn, value: s)}: only changes of \texttt{s}, only
        in that subtree (view decides). \texttt{.animation(\_:)} without
        \texttt{value:} is deprecated — it animated \emph{any} change.
  \item \textbf{Transaction} carries the animation through one update:
        \texttt{withTransaction} is the general form;
        \texttt{.transaction \{ \$0.animation = nil \}} strips an inherited one;
        \texttt{disablesAnimations = true} forces a snap.
  \item \textbf{Animatable}: built-in modifiers interpolate (opacity, offset,
        scale, frame, colour). A custom \texttt{Shape}/modifier must expose
        \texttt{animatableData: VectorArithmetic} (two values $\to$
        \texttt{AnimatablePair}) — otherwise it \emph{jumps}.
  \item \textbf{Transition} (\texttt{.opacity}, \texttt{.move(edge:)}, \\
        \texttt{.asymmetric(insertion:removal:)}) = entering or leaving the tree:
        needs an \textbf{identity change} (\texttt{if}, new \texttt{.id}) \emph{and}
        an animation context, and sits on the conditional view.
  \item \textbf{matchedGeometryEffect(id:in:)}: same id, same
        \texttt{@Namespace}, views mounted \emph{exclusively}, inside
        \texttt{withAnimation} $\to$ hero morph.
  \item \textbf{Springs} (default since iOS 17): \texttt{.spring(duration: 0.5,
        bounce: 0.2)}, \texttt{.smooth}/\texttt{.snappy}/\texttt{.bouncy}.
        Interruptible, re-target \emph{keeping velocity}; curves have fixed duration.
  \item \textbf{iOS 17}: \texttt{PhaseAnimator} steps through discrete phases
        (loop or \texttt{trigger:}); \texttt{KeyframeAnimator} = timed
        per-property \texttt{KeyframeTrack}s (Linear/Spring/Cubic\-Keyframe).
\end{itemize}

\section{Example}
\begin{lstlisting}[language=SwiftSheet]
struct Ring: Shape {
  var p: CGFloat                              // 0...1
  var animatableData: CGFloat { get { p } set { p = newValue } }
  func path(in r: CGRect) -> Path { Path() }  // arc to p
}
struct Card: View {
  @State private var open = false
  @Namespace private var ns
  var body: some View {
    if open { Big().matchedGeometryEffect(id: "c", in: ns) }
    else    { Small().matchedGeometryEffect(id: "c", in: ns) }
    Button("Go") { withAnimation(.snappy) { open.toggle() } }
  }
}
\end{lstlisting}

\section{UIKit}
\begin{itemize}
  \item \texttt{UIView.animate(withDuration:…)}: set \emph{final} values in the
        block. Animatable: \texttt{frame}, \texttt{bounds}, \texttt{center},
        \texttt{transform}, \texttt{alpha}, \texttt{backgroundColor};
        \texttt{text}/\texttt{isHidden} snap. Springs:
        \texttt{usingSpringWithDamping:} (\textless 1 bounces); iOS 17
        \texttt{animate(springDuration:bounce:)}.
  \item \textbf{UIViewPropertyAnimator} is an \emph{object}:
        \texttt{startAnimation()}, \texttt{pauseAnimation()}, scrub
        \texttt{fractionComplete}, \texttt{isReversed},
        finish with \texttt{continueAnimation(withTimingParameters:…)} —
        interruptible gesture-driven UI.
  \item \textbf{Auto Layout}: change \texttt{constant}, then
        \texttt{layoutIfNeeded()} \emph{inside} the block (flush once before).
  \item \texttt{UIView.transition(with:…)} + \texttt{.transitionCrossDissolve};
        custom VC transitions must call \texttt{completeTransition(\_:)}.
\end{itemize}

\columnbreak

\section{Model vs presentation}
\begin{tikzpicture}[sheet, yscale=0.8]
  \draw[->, sheetGrey] (0,0) -- (7,0) node[lbl, below left]{time};
  \draw[->, sheetGrey] (0,0) -- (0,2.4) node[lbl, right]{\texttt{x}};
  \node[lbl, left] at (0,2){100}; \node[lbl, left] at (0,0.08){0};
  \draw[sheetGrey, dashed] (0.8,0) -- (0.8,2.3);
  \node[lbl, below] at (0.8,0){set \texttt{x=100}};
  \draw[sheetGrey, dashed] (5.2,0) -- (5.2,2.3);
  \node[lbl, below] at (5.2,0){done};
  \draw[very thick, sheetBlue] (0,0.08) -- (0.8,0.08) -- (0.8,2) -- (6.8,2);
  \draw[very thick, sheetOrange] (0,0.08) -- (0.8,0.08)
      .. controls (2.0,0.1) and (2.4,2.4) .. (3.3,2.14)
      .. controls (4.1,1.95) and (4.5,2.0) .. (5.2,2) -- (6.8,2);
  \node[lbl, text=sheetBlue, anchor=south west] at (5.3,2.05){\textbf{model}};
  \node[lbl, text=sheetOrange, anchor=west] at (0.95,1.45)
       {\textbf{presentation}};
  \draw[sheetRed, thick, dashed] (2.1,0.95) -- (2.1,2);
  \fill[sheetRed] (2.1,0.95) circle (1.6pt);
  \node[lbl, text=sheetRed, anchor=west, align=left] at (2.25,0.5)
       {drawn: \texttt{presentation()}\\
        touch/\texttt{position}: 100};
\end{tikzpicture}

{\footnotesize UIView animations set the model for you. A raw
\texttt{CABasicAnimation} drives only the presentation $\Rightarrow$ it
\textbf{snaps back} unless you also set the model value.}

\section{When does a transition run?}
\begin{tikzpicture}[sheet]
  \node[st] (c) at (0,0) {state\\change};
  \node[st, draw=sheetGreen, fill=sheetGreen!10] (a) at (1.75,0) {animation\\context?};
  \node[st, draw=sheetGreen, fill=sheetGreen!10] (i) at (3.7,0) {identity\\changed?};
  \node[st, fill=sheetBlue!15] (t) at (5.6,0) {transition\\runs};
  \node[st, draw=sheetRed, fill=sheetRed!10] (s) at (3.7,-1.0) {instant snap};
  \draw[flow] (c) -- (a);
  \draw[flow] (a) -- node[lbl, above]{yes} (i);
  \draw[flow] (i) -- node[lbl, above]{yes} (t);
  \draw[hot] (a) |- node[lbl, pos=0.3, left]{no} (s);
  \draw[hot] (i) -- node[lbl, right]{no: only modifiers animate} (s);
\end{tikzpicture}

\section{Traps}
\begin{itemize}
  \trap{\texttt{withAnimation} animates \emph{all} dependent views — scope
        with \texttt{.animation(\_:value:)} or strip with \texttt{.transaction}.}
  \trap{\texttt{.transition} without an animation = instant; on a stable
        ancestor = never runs.}
  \trap{Custom \texttt{Shape} jumps $\to$ no \texttt{animatableData}.
        matchedGeometry glitch $\to$ both views mounted / two namespaces.}
  \trap{Layer animation snaps back $\to$ model never set
        (\texttt{fillMode = .forwards} + \texttt{isRemovedOnCompletion = false}
        only hides it).}
  \trap{Constraint change jumps $\to$ no \texttt{layoutIfNeeded()} inside the block.}
  \trap{Scrubbing a spring: \texttt{pauseAnimation()} first;
        \texttt{fractionComplete} maps through the curve.}
\end{itemize}

\section{Remember}
\textbf{``State in the block, value in the modifier, identity for
transitions.''} UIKit: \emph{final value in — the model is already there.}
Cheapest to animate: \texttt{transform}, \texttt{opacity}.


\end{multicols}

\noindent{\footnotesize\color{sheetGrey}\textit{Related:} SwiftUI performance \& identity · rendering pipeline · Core Animation (layers, \texttt{CADisplayLink}) · Auto Layout · UIKit touches (custom VC transitions)}

\end{document}
